\documentclass[11pt]{article}

%%%%%%%%%%%%%%%%%%%%%%%%%%%%%%%%%%%%%%%%%%%%%%%%%%%%
% Packages
%%%%%%%%%%%%%%%%%%%%%%%%%%%%%%%%%%%%%%%%%%%%%%%%%%%%

\usepackage{amsmath,amssymb}
\usepackage{graphicx}
\usepackage{booktabs}
\usepackage{algorithm}
\usepackage{algpseudocode}
\usepackage{hyperref}
\usepackage{tikz}
\usepackage{subcaption}

%%%%%%%%%%%%%%%%%%%%%%%%%%%%%%%%%%%%%%%%%%%%%%%%%%%%
\title{SymCorr: A Compiler Framework for Automated Construction and Efficient Evaluation of Lattice QCD Correlation Functions}

\author{...}

\date{}

%%%%%%%%%%%%%%%%%%%%%%%%%%%%%%%%%%%%%%%%%%%%%%%%%%%%

\begin{document}

\maketitle

%%%%%%%%%%%%%%%%%%%%%%%%%%%%%%%%%%%%%%%%%%%%%%%%%%%%
\begin{abstract}

% Motivation

% Challenges of modern lattice spectroscopy

% Need for symbolic automation

% Need for efficient execution

% Main contributions

\end{abstract}

%%%%%%%%%%%%%%%%%%%%%%%%%%%%%%%%%%%%%%%%%%%%%%%%%%%%

\section{Introduction}

\subsection{Motivation}
Lattice QCD provides a non-perturbative framework for studying the properties
and interactions of hadrons directly from quantum chromodynamics. In hadron
spectroscopy, the quantities of interest are commonly extracted from Euclidean
correlation functions constructed using interpolating operators with the
quantum numbers of the states being studied.

Modern spectroscopy calculations often require a basis containing several
interpolating operators rather than a single operator. Different operators may
contain different spin structures, momenta, spatial constructions, or
combinations of hadrons. Correlation matrices constructed from such operator
bases provide the information required to resolve the finite-volume energy
spectrum and to distinguish states with similar quantum numbers.

As the number and complexity of the interpolating operators increase,
constructing the corresponding correlation functions becomes increasingly
cumbersome. For every combination of source and sink operators, the allowed
Wick contractions must be determined, fermionic signs must be treated
consistently, and the flavor, spin, color, and spatial structures must be
carried through to the numerical calculation. Although many of these steps are
systematic, performing them manually for every new operator basis requires
substantial implementation effort and can introduce bookkeeping errors.

The distillation framework~\cite{Peardon:2009gh} provides a particularly useful
setting for spectroscopy calculations because it separates quark propagation
from the operator-dependent structures used to construct hadrons. This makes
it possible to construct a wide range of interpolating operators from a common
set of quark-propagation data. Nevertheless, the transformation from the
mathematical definition of an interpolating operator to the complete numerical
expression for its correlation function remains a non-trivial task.

\textsc{SymCorr} is developed to automate this transformation. Interpolating
operators are first expressed symbolically, after which the package determines
the corresponding Wick contractions and systematically transforms the
resulting expressions toward their numerical representation in the
distillation framework. The aim is to reduce the amount of correlator-specific
algebra and implementation that must be performed manually, while keeping the
physical definition of an operator separate from the details of its numerical
evaluation.

\subsection{Challenges in lattice QCD correlator construction}

For a simple interpolating operator, constructing a two-point correlation
function can be relatively straightforward. A meson operator, for example,
contains one quark and one antiquark field, and the corresponding Wick
contractions involve only a small number of quark lines. The structure becomes
more involved when the operators contain several quark fields or when many
operators are combined into a correlation matrix.

For a basis of interpolating operators \(\{\mathcal O_i\}\), one evaluates the
matrix of correlation functions

\begin{equation}
C_{ij}(t)
=
\left\langle
\mathcal O_i(t)\,\mathcal O_j^{\dagger}(0)
\right\rangle .
\end{equation}

Each source--sink pair can contain different flavor, spin, momentum, and
spatial structures. These structures must be retained consistently while all
allowed contractions between quark and antiquark fields are identified. When
several fields of the same flavor are present, there can be multiple valid
pairings, and the relative signs associated with fermionic permutations must
also be handled correctly.

The amount of bookkeeping therefore grows as one moves from simple meson
operators to baryons and multi-hadron operators. At the same time, realistic
spectroscopy calculations require many source--sink combinations, so the same
basic contraction patterns and operator-dependent structures can appear
repeatedly across different matrix elements.

After the Wick contractions have been determined, the resulting propagator
expressions must still be translated into the objects used for numerical
evaluation. In the distillation framework this means expressing the correlator
in terms of perambulators and operator-dependent structures in the Laplacian
eigenvector space. Carrying out this translation separately for each operator
combination makes the implementation difficult to extend and increases the
amount of repeated algebra that must be checked by hand.

These difficulties motivate a systematic representation of the operators and
their contractions. In \textsc{SymCorr}, the operator definition, Wick
contraction, and subsequent translation to the distillation representation are
treated as distinct stages. This allows each stage to be constructed and
verified independently before introducing the numerical optimizations required
for large correlation matrices.

\subsection{Existing software}
There are several existing software packages that provide tools for constructing 
the interpolating operators in the subduced representations of the lattice symmetry 
groups such as Multi Hadron Interpolators (MHI) \cite{PhysRevD.109.094516} and OpTion\cite{Yan:2025op}. Following this the next step is to perform the Wick contractions and evaluate the resulting expressions. There are several software packages that automate this process and also various algorithms have been proposed to identify the redundancies and optimize the evaluation of the resulting expresssions such as those proposed in \cite{10.1145/3592979.3593409, PhysRevD.110.114505}. However, these packages are not designed to work with the distillation framework and do not provide the necessary abstractions to represent the distillation formalism.
\subsection{Contributions of SymCorr}
In this work we present \textsc{SymCorr}, a domain-spefific compiler framework that automates the construction and optimization of lattice QCD correlation functions within the distillation formalism. The main contributions of this work are as follows:
The principal contributions of this work are

\begin{itemize}

\item A symbolic intermediate representation for lattice QCD interpolating
operators that naturally captures flavor, spin, derivatives, momentum
projection, and color structure.

\item An automatic Wick contraction engine that transforms operator products
into sums of quark-line contraction graphs without manual derivation.

\item A compiler framework that lowers symbolic expressions into optimized
distillation-space execution graphs by identifying reusable computational
structures.

\item A backend capable of generating efficient numerical evaluation routines
that separate symbolic operator construction from architecture-dependent
execution.

\item A modular software design that enables rapid development of new operator
bases while reducing implementation effort and the likelihood of algebraic
errors.

\end{itemize}
%%%%%%%%%%%%%%%%%%%%%%%%%%%%%%%%%%%%%%%%%%%%%%%%%%%%

\section{Overview of SymCorr}

The construction of lattice QCD correlation functions naturally consists of a sequence of symbolic and numerical transformations. Traditionally, these stages are implemented manually, with the symbolic derivation of correlation functions tightly coupled to their numerical realization. As operator bases become larger and more sophisticated, this workflow becomes increasingly difficult to maintain, verify, and extend.

\textsc{SymCorr} addresses this problem by treating correlator construction as a domain-specific compiler. Rather than directly translating interpolating operators into numerical routines, the compiler progressively lowers a high-level symbolic description into an optimized execution plan through a sequence of well-defined compilation stages. This separation of concerns allows symbolic manipulations, algebraic optimizations, and backend-specific execution strategies to be developed independently.

The overall compilation pipeline is illustrated in Fig.~\ref{fig:pipeline}. The process begins with a symbolic representation of interpolating operators, where the fundamental building blocks of lattice QCD---quark fields, gauge links, Dirac matrices, derivatives, and index structures---are represented independently of any numerical backend. The symbolic expressions are subsequently processed by an automatic Wick contraction engine, producing the complete set of fermionic contractions required to evaluate the correlation function.

The contracted expressions are then lowered into a logical execution plan that represents the computation as a collection of backend-independent operations. At this stage, the compiler identifies reusable computational structures and algebraic dependencies while remaining agnostic to the underlying numerical implementation. The logical plan is subsequently transformed into a physical execution plan, where abstract operations are mapped onto concrete numerical primitives appropriate for the chosen backend, such as distillation kernels, perambulators, and tensor contractions.

Finally, the runtime system executes the optimized physical plan on numerical lattice data to produce the desired correlation functions. By separating symbolic operator construction from numerical evaluation, \textsc{SymCorr} provides a modular and extensible framework that simplifies the implementation of new interpolating operators while enabling backend-specific optimizations without modifying the underlying physics description.
%%%%%%%%%%%%%%%%%%%%%%%%%%%%%%%%%%%%%%%%%%%%%%%%%%%%

\section{Symbolic Representation}

The frontend of \textsc{SymCorr} provides a symbolic language for constructing
lattice QCD interpolating operators. Rather than directly generating numerical
tensor contractions, operators are first represented as symbolic expressions
that closely resemble their mathematical definitions. This symbolic
representation preserves the flavor, spin, color, spatial, momentum, and
derivative structure of the operator while remaining completely independent of
the numerical backend. Consequently, algebraic manipulations such as Wick
contractions, expression simplification, and compiler optimizations can be
performed before any numerical implementation is introduced.

Internally, every symbolic expression is represented as an abstract syntax tree
(AST). The leaf nodes correspond to elementary lattice QCD objects such as
quark fields, gamma matrices, gauge links, derivatives, momentum projections,
and tensor objects, while internal nodes represent algebraic operations
including multiplication, summation, tensor contraction, and operator
composition. This tree-based representation preserves the complete symbolic
structure of the interpolating operator and forms the intermediate
representation on which subsequent compiler passes operate.

In its most general form, an interpolating operator may be regarded as an
algebraic composition of elementary symbolic objects,

\begin{equation}
\mathcal{O}
=
\mathcal{C}
\left[
\prod_{i=1}^{N}
\Phi_i
\right],
\label{eq:generic_operator}
\end{equation}

where

\begin{equation}
\Phi_i
\in
\left\{
q,\,
\bar q,\,
\Gamma,\,
U,\,
D,\,
P,\,
\epsilon,\,
\delta,\,
\ldots
\right\},
\end{equation}

and $\mathcal C$ denotes the complete contraction of color, spin, flavor,
spatial, and other auxiliary indices required to project the operator onto the
desired quantum numbers. The elementary symbolic objects represented by
\textsc{SymCorr} are described below.

\subsection{Quark fields}
The fundamental fermionic objects are quark and antiquark fields,

\begin{equation}
q_f^{a\alpha}(x),
\qquad
\bar q_f^{a\alpha}(x),
\end{equation}

where $f$ denotes the flavor index, $a$ the color index, $\alpha$ the Dirac
spin index, and $x$ the spacetime coordinate. These symbolic fields constitute
the elementary degrees of freedom from which all interpolating operators are
constructed.

Since the frontend manipulates symbolic objects rather than numerical fields,
the representation is independent of the underlying fermion discretization,
allowing the same symbolic description to be compiled for different numerical
backends.

\subsection{Indices}

Every symbolic object carries a collection of indices describing its
transformation properties. Depending on the object, these include color, spin,
flavor, spacetime, momentum, derivative, and displacement indices. Throughout
the compilation process these indices remain symbolic, allowing the compiler to
perform consistency checks, identify valid tensor contractions, and preserve
the algebraic structure of the operator until numerical code generation.

\subsection{Dirac matrices}

Dirac matrices are represented as symbolic operators acting on spinor indices.
They may be freely combined with quark fields to construct arbitrary fermion
bilinears and higher-order interpolating operators. Treating Dirac matrices
symbolically enables the compiler to manipulate operator expressions without
prematurely expanding the associated spin algebra.

\subsection{Gauge transporters}

Gauge transporters are represented by symbolic Wilson lines,

\begin{equation}
U(P)
=
\prod_{\ell\in P}
U_\ell,
\end{equation}

where $P$ denotes an arbitrary lattice path. These objects ensure gauge
covariance for non-local interpolating operators and displaced quark fields.
Because the path itself is represented symbolically, arbitrary displacement
patterns can be incorporated without modifying the compiler.

\subsection{Covariant derivatives}

Covariant derivatives are represented as symbolic operators acting on quark
fields. They may be combined to construct arbitrary derivative structures and
operators carrying non-zero orbital angular momentum. Their symbolic
representation preserves the derivative structure throughout the compilation
pipeline before being lowered to the corresponding backend-specific
implementation.

\subsection{Momentum projections}

Momentum projection operators are represented by symbolic phase factors,

\begin{equation}
P_{\mathbf p}(x)
=
e^{i\mathbf p\cdot\mathbf x},
\end{equation}

which may act on arbitrary symbolic expressions. By retaining momentum
projections symbolically, \textsc{SymCorr} preserves translational properties
and momentum conservation until the numerical execution stage.

\subsection{Tensor objects}

Many interpolating operators require additional tensors to contract color,
spin, flavor, or lattice symmetry indices. These include the Kronecker delta,
the Levi--Civita tensor,

\begin{equation}
\delta_{ij},
\qquad
\epsilon_{abc},
\end{equation}

as well as arbitrary user-defined symbolic tensors. These objects are treated
uniformly within the symbolic representation and participate in the same
algebraic manipulation and contraction framework as the remaining symbolic
objects.

\subsection{Interpolating operators}

Interpolating operators are obtained by composing the elementary symbolic
objects introduced above through algebraic operations represented in the AST.
The resulting expressions may describe local or non-local operators, mesons,
baryons, tetraquarks, pentaquarks, hybrid operators, or arbitrary multi-hadron
systems.

For example, a generic meson interpolating operator may be written as

\begin{equation}
\mathcal O_M(x)
=
\bar q_f(x)\,
\Gamma\,
D^{(n)}
U(P)\,
P_{\mathbf p}(x)\,
q_g(x),
\end{equation}

whereas a baryon interpolating operator takes the form

\begin{equation}
\mathcal O_B(x)
=
\epsilon_{abc}
\left(
q_1^a(x)
\Gamma_1
q_2^b(x)
\right)
\Gamma_2
q_3^c(x).
\end{equation}

More complicated interpolators are constructed recursively by composing these
elementary building blocks. Since the complete symbolic structure is preserved,
the resulting expressions can subsequently undergo automatic Wick contraction,
compiler optimization, and backend-specific lowering without requiring any
manual algebraic manipulation.

%%%%%%%%%%%%%%%%%%%%%%%%%%%%%%%%%%%%%%%%%%%%%%%%%%%%

\section{Automatic Wick Contraction}

Following the symbolic construction of the interpolating operators,
\textsc{SymCorr} automatically applies Wick's theorem to generate the
corresponding correlation function. This stage transforms products of fermionic
field operators into sums over quark propagator contractions while preserving
the complete symbolic structure of the expression.

Unlike traditional workflows, where Wick contractions are derived manually for
each new operator basis, the symbolic representation allows the contraction
process to be fully automated. The output of this stage is a symbolic
expression containing all valid contraction terms together with the associated
fermionic permutation signs. These expressions form the input to the
optimization stages described in the following sections.

\subsection{Wick's theorem}

Consider a generic Euclidean correlation function

\begin{equation}
C(t)
=
\left<
\mathcal O_{\mathrm{sink}}(t)
\mathcal O_{\mathrm{source}}^\dagger(0)
\right>.
\end{equation}

Applying Wick's theorem replaces every quark--antiquark pair by the
corresponding propagator,

\begin{equation}
q(x)\bar q(y)
\rightarrow
S(x,y),
\end{equation}

and expresses the correlation function as a sum over all complete
fermionic contractions. Since QCD conserves flavor, contractions are permitted
only between quark and antiquark fields carrying identical flavor quantum
numbers.

\subsection{Contraction graph generation}

The contraction engine first identifies every quark and antiquark field
appearing in the symbolic expression. For each quark flavor, all complete
matchings between quark and antiquark fields are generated subject to flavor
conservation. Every valid matching corresponds to one contraction graph.

If a correlation function contains \(n_f\) quark fields of flavor \(f\), the
number of possible flavor-preserving contractions is

\begin{equation}
N_{\mathrm{contr}}
=
\prod_f n_f!.
\end{equation}

The symbolic representation enables these contraction graphs to be generated
independently of any numerical backend, allowing subsequent compiler passes to
operate directly on the resulting algebraic structure.

\subsection{Fermionic permutations}
The contraction of Grassmann-valued fermion fields introduces a sign determined
by the parity of the permutation required to bring contracted fields adjacent
to one another. Consequently, every contraction graph is associated with a
fermionic sign

\begin{equation}
(-1)^{P},
\end{equation}

where \(P\) denotes the parity of the corresponding permutation.

During contraction generation, \textsc{SymCorr} explicitly tracks these
permutations so that the correct fermionic phase accompanies every contraction
term throughout the compilation pipeline.

\subsection{Determinant representation}

The antisymmetric nature of Grassmann fields allows many contraction
expressions to be written compactly as determinants of quark propagators.
Rather than explicitly enumerating every permutation, groups of contractions
associated with identical fermion species can be represented by determinant
expressions, leading to a substantial reduction in symbolic complexity.

Recent work~\cite{PhysRevD.110.114505} demonstrated that determinant
representations together with polynomial identity testing can significantly
reduce the number of independent determinants required for conventional lattice
QCD correlators. Although this construction is highly effective for point
propagators, the distillation framework introduces additional mode-space kernel
structures that prevent its direct application.

In \textsc{SymCorr}, the determinant formulation is retained as the underlying
symbolic representation, while additional compiler optimizations are introduced
to identify and eliminate redundant computations arising specifically within
the distillation formalism. These optimizations are discussed in
Section~\ref{sec:optimization}.

\subsection{Expression simplification}

The contraction stage generally produces expressions containing algebraically
equivalent terms and redundant symbolic structures. Before lowering the
expressions into the compiler intermediate representation, \textsc{SymCorr}
applies a collection of symbolic simplification rules.

These include the elimination of identical contraction terms, simplification
of tensor contractions, canonical ordering of symbolic objects, and the
application of algebraic identities involving Dirac matrices, color tensors,
and momentum projections. Performing these simplifications symbolically
reduces the size of the generated execution graph and improves the efficiency
of subsequent optimization and numerical evaluation stages.

%%%%%%%%%%%%%%%%%%%%%%%%%%%%%%%%%%%%%%%%%%%%%%%%%%%%

\section{Compilation to the Distillation Representation}

The symbolic expressions produced by the automatic Wick contraction stage are
written in terms of quark propagators connecting contracted quark and
antiquark fields. While these expressions are mathematically complete, they are
not directly suitable for efficient numerical evaluation within the
distillation framework. Instead, the quark propagators must be rewritten in
terms of the elementary objects of distillation, namely the perambulators and
mode-space operator kernels.

The purpose of this compilation stage is therefore to lower the symbolic Wick
expressions into a backend-independent intermediate representation expressed
entirely in terms of distillation primitives. During this transformation,
\textsc{SymCorr} identifies reusable operator kernels, factors propagators into
their distillation representation, and constructs an execution graph that
explicitly describes the numerical operations required to evaluate the
correlation function. Since the compilation is performed symbolically, the
resulting representation remains independent of the numerical backend while
exposing opportunities for subsequent optimization and task scheduling.

The compilation pipeline implemented in this stage consists of the following
steps:

\begin{enumerate}
    \item Rewrite every quark propagator using the distillation decomposition.
    \item Identify maximal mode-space operator kernels.
    \item Construct reusable kernel descriptors.
    \item Generate backend-independent contraction tasks.
    \item Assemble the execution graph for numerical evaluation.
\end{enumerate}

The following subsections describe each of these stages in detail.
\subsection{Distillation formalism}
Distillation introduces a low-rank approximation of the quark fields by projecting them onto a subspace spanned by the lowest eigenmodes of the gauge-covariant Laplacian operator. Where the laplacian operator is defined as
\begin{equation}
\Delta = \sum_{\mu} \left( 2 - U_\mu(x) - U_\mu^\dagger(x-\hat{\mu}) \right),
\end{equation}
and the eigenvectors \(v_k(x)\) satisfy
\begin{equation}
\Delta v_k(x) = \lambda_k v_k(x),
\end{equation}
where \(\lambda_k\) are the corresponding eigenvalues. The distillation operator is then
\begin{equation}
\Box = \sum_{k=1}^{N_v} v_k(x) v_k^\dagger(y),
\end{equation}
where \(N_v\) is the number of eigenvectors retained in the projection. The smeared quark fields are defined as
\begin{equation}
\tilde{q}(x) = \Box q(x), \quad \tilde{\bar{q}}(x) = \bar{q}(x) \Box.
\end{equation}

\subsection{Perambulators}
The quark propagators in the distillation framework are replaced by perambulators, which are defined as
\begin{equation}
\tau_{kl}(t, t') = v_k^\dagger(t) S(t, t') v_l(t'),
\end{equation}
where \(S(t, t')\) is the quark propagator between timeslices \(t\) and \(t'\), and \(v_k(t)\) are the eigenvectors of the Laplacian operator at timeslice \(t\). The perambulators encode the propagation of the smeared quark fields between different timeslices and are essential for constructing correlation functions in the distillation formalism.They are typically precomputed and stored for reuse across multiple correlation function evaluations, significantly reducing the computational cost of the overall calculation.
\subsection{Kernel identification}
The distillation compilation stage identifies the mode-space operator kernels that arise from the contraction expressions. These kernels are defined as
\begin{equation}
K_{kl}(t) = v_k^\dagger(t) \Gamma D^{(n)} U(P) P_{\mathbf p}(t) v_l(t),
\end{equation}
where \(\Gamma\) is a Dirac matrix, \(D^{(n)}\) is a covariant derivative operator, \(U(P)\) is a gauge transporter along path \(P\), and \(P_{\mathbf p}(t)\) is a momentum projection operator.
The kernels capture the action of the interpolating operators in the mode space defined by the Laplacian eigenvectors. By identifying and reusing these kernels across different contraction terms, the compiler can reduce the number of unique computations required for evaluating the correlation functions, leading to improved performance and reduced memory usage. 
\subsection{Kernel recipes}

\subsection{Contraction tasks}

%%%%%%%%%%%%%%%%%%%%%%%%%%%%%%%%%%%%%%%%%%%%%%%%%%%%

\section{Compiler Backend}

%%%%%%%%%%%%%%%%%%%%%%%%%%%%%%%%%%%%%%%%%%%%%%%%%%%%

\subsection{Logical Execution Plan}

Explain

LogicalTask

LogicalExecutionPlan

%%%%%%%%%%%%%%%%%%%%%%%%%%%%%%%%%%%%%%%%%%%%%%%%%%%%

\subsection{Static Analyses}

DependencyAnalysis

LifetimeAnalysis

%%%%%%%%%%%%%%%%%%%%%%%%%%%%%%%%%%%%%%%%%%%%%%%%%%%%

\subsection{Physical Lowering}

Logical Plan

↓

Physical Execution Plan

%%%%%%%%%%%%%%%%%%%%%%%%%%%%%%%%%%%%%%%%%%%%%%%%%%%%

\subsection{Physical Instruction Set}

Allocate

Load

Evaluate

Contract

Copy

Store

Free

Barrier

Wait

Event

Prefetch

%%%%%%%%%%%%%%%%%%%%%%%%%%%%%%%%%%%%%%%%%%%%%%%%%%%%

\subsection{Plan Verification}

%%%%%%%%%%%%%%%%%%%%%%%%%%%%%%%%%%%%%%%%%%%%%%%%%%%%

\section{Optimization Framework}

%%%%%%%%%%%%%%%%%%%%%%%%%%%%%%%%%%%%%%%%%%%%%%%%%%%%

\subsection{Planning Pipeline}

%%%%%%%%%%%%%%%%%%%%%%%%%%%%%%%%%%%%%%%%%%%%%%%%%%%%

\subsection{Kernel Reuse}

%%%%%%%%%%%%%%%%%%%%%%%%%%%%%%%%%%%%%%%%%%%%%%%%%%%%

\subsection{Memory Planning}

%%%%%%%%%%%%%%%%%%%%%%%%%%%%%%%%%%%%%%%%%%%%%%%%%%%%

\subsection{Analytical Cost Model}

%%%%%%%%%%%%%%%%%%%%%%%%%%%%%%%%%%%%%%%%%%%%%%%%%%%%

\section{Runtime System}

%%%%%%%%%%%%%%%%%%%%%%%%%%%%%%%%%%%%%%%%%%%%%%%%%%%%

\subsection{Execution Context}

%%%%%%%%%%%%%%%%%%%%%%%%%%%%%%%%%%%%%%%%%%%%%%%%%%%%

\subsection{Resource Manager}

%%%%%%%%%%%%%%%%%%%%%%%%%%%%%%%%%%%%%%%%%%%%%%%%%%%%

\subsection{Storage Manager}

%%%%%%%%%%%%%%%%%%%%%%%%%%%%%%%%%%%%%%%%%%%%%%%%%%%%

\subsection{Kernel Evaluation}

%%%%%%%%%%%%%%%%%%%%%%%%%%%%%%%%%%%%%%%%%%%%%%%%%%%%

\subsection{Task Execution}

%%%%%%%%%%%%%%%%%%%%%%%%%%%%%%%%%%%%%%%%%%%%%%%%%%%%

\subsection{Output Assembly}

%%%%%%%%%%%%%%%%%%%%%%%%%%%%%%%%%%%%%%%%%%%%%%%%%%%%

\section{Worked Examples}

%%%%%%%%%%%%%%%%%%%%%%%%%%%%%%%%%%%%%%%%%%%%%%%%%%%%

\subsection{Meson correlator}

%%%%%%%%%%%%%%%%%%%%%%%%%%%%%%%%%%%%%%%%%%%%%%%%%%%%

\subsection{Baryon correlator}
%%%%%%%%%%%%%%%%%%%%%%%%%%%%%%%%%%%%%%%%%%%%%%%%%%%%

\subsection{Multi-hadron correlator}

%%%%%%%%%%%%%%%%%%%%%%%%%%%%%%%%%%%%%%%%%%%%%%%%%%%%

\section{Performance Evaluation}

%%%%%%%%%%%%%%%%%%%%%%%%%%%%%%%%%%%%%%%%%%%%%%%%%%%%

\subsection{Compiler performance}

Compilation time

%%%%%%%%%%%%%%%%%%%%%%%%%%%%%%%%%%%%%%%%%%%%%%%%%%%%

\subsection{Runtime performance}

Execution time

%%%%%%%%%%%%%%%%%%%%%%%%%%%%%%%%%%%%%%%%%%%%%%%%%%%%

\subsection{Kernel reuse}

%%%%%%%%%%%%%%%%%%%%%%%%%%%%%%%%%%%%%%%%%%%%%%%%%%%%

\subsection{Memory reduction}

%%%%%%%%%%%%%%%%%%%%%%%%%%%%%%%%%%%%%%%%%%%%%%%%%%%%

\subsection{Scalability}

%%%%%%%%%%%%%%%%%%%%%%%%%%%%%%%%%%%%%%%%%%%%%%%%%%%%

\section{Applications}

Discuss

Meson spectroscopy

Baryons

Multi-hadron operators

Coupled-channel calculations

Future large operator bases

%%%%%%%%%%%%%%%%%%%%%%%%%%%%%%%%%%%%%%%%%%%%%%%%%%%%

\section{Future Extensions}

Profile-guided optimisation

Multi-GPU

Distributed execution

Topology-aware scheduling

Additional optimisation passes

%%%%%%%%%%%%%%%%%%%%%%%%%%%%%%%%%%%%%%%%%%%%%%%%%%%%

\section{Conclusions}

%%%%%%%%%%%%%%%%%%%%%%%%%%%%%%%%%%%%%%%%%%%%%%%%%%%%

\bibliographystyle{unsrt}
\bibliography{references}

\end{document}